\scanpage{63}
\mspage{56}

est dans $\mathbb{W}$, on veut que le foncteur $u^*$ induise une \uline{équivalence}

\[ u^{*}_{\mathrm{l.c.}}:\mathrm{HO}^{\unsure{\mathrm{lc}|\ell\mathrm{c}}}_W(I)\xrightarrow{\approx}\mathrm{HO}^{\mathrm{lc}}_W(I'). \]

Mais je doute qu'on soit armé pour vérifier un axiome aussi fort, sans avoir \inserted{au moins} une théorie des opérations $u_!$ ou $u_*$.

On pourrait à \deleted{tout} hasard poser cet axiome \deleted{\unsure{pour}} en prenant $\mathbb{W}=\mathbb{W}_{\unsure{\alpha|a}}\overset{\text{déf}}{\unsure{\ne|=}}$ le plus petit loc. fond satisfaisant $W(1\text{ à }7)$ (le plus grand de ces localiseurs étant $\mathbb{W}_\infty$). Pour l'instant, la façon de vérifier cet axiome \deleted{\unsure{(qu}} (en termes de la donnée $(M,W)$) est pour moi un mystère $\pm$ total. (Mais \unsure{le|je} sais le vérifier pour $\mathrm{HOT}_W$, quand $W\subset\mathrm{Fl}(\mathrm{Cat})$ satisfait $W(1,2,3,4,5\ldots)$

Je vais m'en occuper, une fois que j'aurai une maîtrise raisonnable sur le formalisme de $u_*,u_!$ relatif \unsure{aux|au} $u:I'\to I$ dans $\mathrm{Diag}$. L'hypothèse quillénisable \inserted{seule} semble trop grossière pour assurer l'existence de $u_{\unsure{!|*}}$ ou de $u_*$ dans des cas intéressants.
% archive sheet number: 62
